\begin{abstract}
We study how stablecoins affect parallel foreign exchange markets in economies with fixed exchange rate regimes. We show that stablecoins are not merely a cheaper transactional instrument, but also act as an endogenous public signal that aggregates dispersed information about exchange rate misalignment. Embedding this mechanism in a global-games framework, we demonstrate that greater stablecoin usage increases the precision of public information, compresses belief dispersion, and strengthens coordination among agents. 

This generates a fundamental tradeoff: while stablecoins improve efficiency and risk-sharing at low levels of misalignment, they increase fragility by facilitating coordinated speculative attacks at higher levels. Quantitatively, stablecoin demand rises with overvaluation, crisis probabilities increase, and welfare is non-monotonic. Our results highlight how financial innovations that improve price discovery can destabilize equilibrium outcomes in environments with strategic complementarities.
\end{abstract}